% IEEE CloudNet 2026 — Demo & Artifact Track
% 2 pages, 10 pt, double-column
% Compile: pdflatex main && pdflatex main
%
\documentclass[10pt,conference]{IEEEtran}

\usepackage[T1]{fontenc}
\usepackage[utf8]{inputenc}
\usepackage{booktabs}
\usepackage{url}
\usepackage{cite}
\usepackage{makecell}
\usepackage[table]{xcolor}
\usepackage{amssymb}

% ─────────────────────────────────────────────────────────────────────────────
\title{QoS-Aware AI Inference Offloading for Vehicular Edge Computing:\\
       A Physics-Based Dataset and Simulation Framework}

\author{%
  \IEEEauthorblockN{%
    Konstantinos Gyftodimos,\quad
    Eirini Liotou,\quad
    Kyriakos Chiotis,\quad
    George Dimitrakopoulos%
  }%
  \IEEEauthorblockA{%
    Dept.\ of Informatics and Telematics,
    Harokopio University of Athens (HUA), Athens, Greece\\
    \{kgyftodimos,\;eliotou,\;kchiotis,\;gdimitra\}@\textit{hua.gr}\quad
    \textit{Presenter: K.\,Gyftodimos\,---\,onsite}%
  }%
}

% ─────────────────────────────────────────────────────────────────────────────
\begin{document}
\maketitle

% ── Abstract ─────────────────────────────────────────────────────────────────
\begin{abstract}
We present a physics-based open-source simulation framework for QoS-aware AI
inference offloading in vehicular edge computing, comprising a 50,000-sample
synthetic dataset across five road scenarios, a multi-output machine-learning
predictor, and a real-time closed-loop demonstration.
\end{abstract}

% ── I. Overview ──────────────────────────────────────────────────────────────
\section{Overview}

Autonomous vehicles generate continuous high-rate video streams that must be
processed by YOLOv8~\cite{yolov8} object detection under a strict 100\,ms
end-to-end latency constraint.  The simulation framework models a four-tier
execution hierarchy: (A)~a dedicated lightweight on-board neural processing
unit (NPU), (B)~the vehicle's primary high-throughput GPU, (C)~a nearby
Multi-Access Edge Computing (MEC)~\cite{etsi_mec} server, and (D)~a remote
cloud data centre.  Each tier runs a different YOLOv8 variant and offers a
distinct latency--accuracy trade-off (Table~\ref{tab:modes}).  Selecting the
optimal tier in real time requires jointly predicting network conditions and
enforcing Quality of Service (QoS) constraints simultaneously.

Three properties of the vehicular setting make that decision hard.  The radio
link is non-stationary: mobility, NLOS blockage, and cell-edge handovers
change the achievable uplink rate on sub-second timescales, so a decision
taken on stale state can miss the deadline.  The MEC server is contended, so
its queueing delay depends on neighbouring vehicles the ego-vehicle cannot
observe.  And the tiers differ in accuracy as well as in latency, so
protecting the latency budget can silently degrade perception---a safety
concern, not merely a performance one.

Most vehicular offloading studies rely on proprietary radio traces, or
abstract the channel so far that MEC latency variability is not faithfully
reproduced.  This artifact provides an open, reproducible baseline---from
3GPP-calibrated channel simulation to closed-loop optimisation---that is
self-contained, runs offline, and ports to any compute-intensive edge
inference workload.

\textbf{Delivery format:} onsite demonstration; the system runs as a single
terminal command on a standard laptop, with no external hardware or network.

% ── II. Innovation \& Contribution ───────────────────────────────────────────
\section{Innovation \& Contribution}

\setlength{\itemsep}{0pt}\setlength{\parskip}{0pt}
\begin{enumerate}
  \setlength{\itemsep}{2pt}
  \item \textbf{Physics-based synthetic dataset (50\,k samples).}
        Generated via a 3GPP TR\,38.901-calibrated log-distance path-loss
        model~\cite{3gpp38901} with log-normal shadowing and non-line-of-sight
        (NLOS) transitions, combined with Shannon--Hartley
        throughput~\cite{shannon1948}.  Five road scenarios span 300\,m
        to 10\,km MEC inter-site distances (urban dense through rural),
        with scenario-specific speed, bandwidth, and frequency profiles.
        \emph{How the data is acquired:} no proprietary or measured trace is
        used---every row is synthesised.  A scenario is drawn from a fixed
        prior, then speed, MEC distance, co-served UEs (Poisson, rush-hour
        scaled), server load (Beta), payload, and WAN congestion are sampled
        from scenario-conditioned distributions and pushed through a
        deterministic chain: path loss $\to$ RSSI (log-normal shadowing,
        25--45\,dB NLOS penalty) $\to$ SINR $\to$ Shannon rate.  RTT labels
        sum uplink, propagation, load-dependent queueing, and a 20\,ms
        handover penalty; reachability is labelled by an SINR threshold.
        Labels are generated, not annotated, and the table regenerates from
        seed~42 in under a minute (MIT license; schema in Sec.~\ref{sec:artifact}).

  \item \textbf{Multi-output machine-learning (ML) predictor.}
        A \texttt{RandomForestRegressor} (200 trees, 80/20 train-test split)
        predicts MEC and cloud round-trip time (RTT) and server capacity
        utilisation; a \texttt{GradientBoostingClassifier} predicts MEC
        availability (binary).  All scikit-learn~\cite{sklearn} pipelines are
        serialised to disk after training and loaded on subsequent runs,
        reducing cold-start time to $\approx$1\,s.  Non-stationarity is
        handled at decision time rather than by the regressor alone: an
        exponential moving average ($\alpha$\,=\,0.3) damps jitter, while
        spikes the model cannot anticipate---abrupt NLOS blockage, a
        handover---are absorbed by the closed loop, which re-evaluates every
        frame at 5--10\,Hz, so a mispredicted frame costs one frame interval
        rather than a persistent mode error.  Reachability is a separate
        classifier, so a hard outage is a discrete event and not an
        extrapolated latency, and a shipped 90th-percentile model
        (\texttt{predictor\_q90.pkl}) can replace the mean predictor for
        conservative gating.  Tracking genuine drift would still need online
        re-fitting; the training entry point is exposed for that.

  \item \textbf{Edge-first QoS optimizer.}
        The decision engine enforces three hard QoS budgets: total latency
        $\leq$100\,ms, accuracy loss $\leq$15\%, and server capacity
        $<$95\%.  Among feasible modes it applies an edge-first preference
        rule: Mode~C (MEC) is chosen over Mode~D (Cloud) when the cloud
        model's incremental mean Average Precision (mAP) gain is below
        5\,pts, conserving Wide Area Network (WAN) bandwidth.
        If no mode satisfies all constraints---e.g.\ during a connectivity
        outage---the optimizer falls back to Mode~B (on-board GPU) with an
        explicit reason flag (\texttt{fallback\_no\_feasible\_mode}).
        Mode~B is deliberately allowed to break the accuracy budget: the
        budgets are ranked lexicographically, with latency and reachability
        held as hard safety constraints and accuracy loss the one that yields
        when nothing is feasible---through a prolonged outage, degraded
        detection beats none.  The violation is never silent: the flag is
        raised on every affected frame, the signal a production stack forwards
        to the planner to reduce speed or extend headway until a compliant tier
        returns.

  \item \textbf{Live four-tier failover demo.}
        One command exercises all four tiers in a 100-frame run: an injected
        outage (frames 40--54) drives a full C\,$\to$\,B\,$\to$\,C failover,
        making every QoS constraint path observable in real time.
\end{enumerate}

% Merged spec + live snapshot table; floats to the top of a page-2 column.
\begin{table}[t]
\centering
\caption{%
  YOLOv8 execution modes: specification and live optimizer snapshot
  (suburban, frame 100/100, MEC dist.\ 611\,m, signal $-$55.5\,dBm,
  4 mode switches).  Tot.\,(ms)\,=\,inference\,$+$\,network RTT.%
}
\label{tab:modes}
\scriptsize\setlength{\tabcolsep}{3pt}
\begin{tabular}{@{}clrrrl@{}}
\toprule
\textbf{M.} & \textbf{Hardware\,/\,Model} &
  \textbf{Tot.\,(ms)} & \textbf{mAP} & \textbf{AccLoss} &
  \textbf{Status} \\
\midrule
A & Car NPU\,/\,\texttt{yolov8n} &  8.0 & 37.3 & 20.8\% &
  $\times$\ loss\,$>$\,15\% \\
B & Car GPU\,/\,\texttt{yolov8n} &  5.0 & 37.3 & 20.8\% &
  $\times$\ loss\,$>$\,15\% \\
\rowcolor{gray!15}
C & MEC\,/\,\texttt{yolov8m}     & 56.9 & 50.2 &  4.9\% &
  $\blacktriangleleft$\ \textbf{SELECTED} \\
D & Cloud\,/\,\texttt{yolov8l}   & 71.6 & 52.9 &  0.0\% &
  feasible, not preferred \\
\bottomrule
\multicolumn{6}{@{}l}{Budget: latency $\leq$100\,ms, acc.\ loss $\leq$15\%, cap.\ $<$95\%.}\\
\multicolumn{6}{@{}l}{Modes A/B exceed the accuracy budget; emergency fallback only.}
\end{tabular}
\end{table}

% ── III. Relevance to CloudNet ────────────────────────────────────────────────
\section{Relevance to CloudNet}

The artifact targets CloudNet researchers working on vehicular edge
computing~\cite{shi2016edge}, task offloading~\cite{chen2019offloading}, and
5G edge AI~\cite{liu2020edge}.  It contributes a \textbf{reproducible
baseline}---a 3GPP-calibrated offloading dataset over five deployment
archetypes, from dense urban (300\,m MEC spacing) to rural (10\,km)---and
\textbf{evidence that offloading is necessary}: on-board YOLOv8n loses
20.8\% accuracy against cloud YOLOv8l, breaching the 15\% budget whenever a
remote tier is reachable.  It is also an \textbf{end-to-end blueprint}
coupling ML network-state prediction to a policy-based QoS optimizer,
portable to other workloads (LiDAR, V2X perception) and extendable with new
tiers or policies.

% ── IV. Demonstration Plan \& Requirements ────────────────────────────────────
\section{Demonstration Plan \& Requirements}

\textbf{Setup.}
Run \texttt{pip install -r requirements.txt}, then \texttt{python demo.py}.
The first run auto-generates the 50\,k dataset and trains all models
($\approx$50\,s); later runs load the serialised models and start in
$\approx$2\,s.

\textbf{What the audience observes.}
The terminal view (Table~\ref{tab:modes}) refreshes at 5\,Hz for 100 frames,
showing all four tiers at once with live per-mode latency, accuracy loss,
capacity, and decision reason.  In normal suburban conditions Mode~C is
selected over Mode~D by the edge-first rule, with Modes~A and~B visible and
marked infeasible.  At frame~40 the injected outage forces MEC unavailable and
cloud RTT to 200\,ms; the optimizer drops to Mode~B with reason
\texttt{fallback\_no\_feasible\_mode}, then recovers.

\textbf{AV / Power:} standard laptop (CPU-only), terminal with ANSI colour;
no special power, networking, or AV equipment needed.

% ── V. Artifact Details ───────────────────────────────────────────────────────
\section{Artifact Details}
\label{sec:artifact}

All components are in a public GitHub repository
(\url{https://github.com/icsa-hua/CLOUDNET2026-QoS-Offloading}, release
\texttt{v1.0.0}; Ubuntu~22.04 / macOS~14 / Python~3.11).  The
\textbf{dataset} is a 50\,k-row CSV regenerated by
\texttt{python -m src.dataset.generate -{}-seed 42}.  Its 13 observable
features cover the scenario label, vehicle speed, time of day, MEC distance,
RSSI, channel bandwidth, MEC load, co-served UEs, handover flag, payload size,
WAN congestion, cloud base RTT, and cloud load; the 5 targets are MEC
reachability (binary), MEC and cloud RTT (ms), and MEC and cloud capacity
(\%).  The
\textbf{framework} splits into \texttt{src/dataset/} (channel sim),
\texttt{src/predictor/} (ML pipelines), and \texttt{src/offloading/} (QoS
optimizer); dataset and models are bundled and auto-recreated if absent.  The
\textbf{demo} (\texttt{python demo.py}) is configurable via
\texttt{-{}-scenario}, \texttt{-{}-hz}, and \texttt{-{}-outage-at}.

% ── VI. Ethics, Privacy \& Safety ────────────────────────────────────────────
\section{Ethics, Privacy \& Safety}

The dataset is entirely synthetic---generated from a calibrated radio-physics
model with no real-world measurements, user data, or personally identifiable
information.  No human subjects or consent procedures are involved; the
framework poses no safety risk.

% ── Acknowledgment ───────────────────────────────────────────────────────────
\section*{Acknowledgment}

This work has been funded by the Shift2SDV project (Grant Agreement
No.~101194245), supported by the Chips Joint Undertaking and its members,
including top-up funding by the national authorities of Austria, Denmark,
Germany, Greece, Finland, Italy, Netherlands, Poland, Portugal, Spain, and
Turkey.

\begin{thebibliography}{8}
\setlength{\itemsep}{0pt plus 0.3pt}

\bibitem{yolov8}
G.~Jocher, A.~Chaurasia, and J.~Qiu,
``Ultralytics YOLOv8,'' 2023.
\url{https://github.com/ultralytics/ultralytics}

\bibitem{etsi_mec}
ETSI,
``GS MEC~003: MEC framework and reference architecture, v3.1.1,'' 2022.

\bibitem{3gpp38901}
3GPP,
``TR~38.901: Channel model for frequencies from 0.5 to 100~GHz,
v16.1.0,'' 2019.

\bibitem{shannon1948}
C.~E.~Shannon,
``A mathematical theory of communication,''
\textit{Bell Syst. Tech. J.}, vol.~27, pp.~379--423, 1948.


\bibitem{sklearn}
F.~Pedregosa \textit{et~al.},
``Scikit-learn: Machine learning in Python,''
\textit{J. Mach. Learn. Res.}, vol.~12, pp.~2825--2830, 2011.

\bibitem{shi2016edge}
W.~Shi, J.~Cao, Q.~Zhang, Y.~Li, and L.~Xu,
``Edge computing: Vision and challenges,''
\textit{IEEE Internet Things J.}, vol.~3, no.~5, pp.~637--646, 2016.

\bibitem{chen2019offloading}
X.~Chen \textit{et~al.},
``Optimized computation offloading via deep reinforcement learning,''
\textit{IEEE Internet Things J.}, vol.~6, no.~3, pp.~4005--4018, 2019.

\bibitem{liu2020edge}
Y.~Liu \textit{et~al.},
``Toward edge intelligence: MEC for 5G and IoT,''
\textit{IEEE Internet Things J.}, vol.~7, no.~8, pp.~6722--6747, 2020.

\end{thebibliography}

\end{document}
